% Fig. 2b - Escenario de colisión en Legacy
% Compile: pdflatex fig2_legacy_collision.tex

\documentclass[border=5pt]{standalone}

\usepackage{amsmath}
\usepackage{tikz}
\usetikzlibrary{arrows.meta, positioning, calc}

\begin{document}

\begin{tikzpicture}[
    >=Stealth,
    font=\small,
    entity/.style={
        draw,
        rectangle,
        minimum width=2cm,
        minimum height=0.55cm,
        fill=black!5,
        font=\small\bfseries,
        thick
    },
    msg/.style={
        ->,
        semithick
    },
    step/.style={
        font=\small\bfseries\itshape,
        text=black!65,
        anchor=east
    },
    sidenote/.style={
        font=\small\itshape,
        text=black!70,
        align=left
    },
    callout/.style={
        draw=black!40,
        fill=black!3,
        rounded corners=2pt,
        inner sep=4pt,
        font=\small,
        align=left,
        line width=0.4pt,
        text width=5cm
    }
]

% Entidades
\node[entity] (ue1) at (0,0) {UE$_1$};
\node[entity] (ue2) at (2.5,0) {UE$_2$};
\node[entity] (enb) at (8,0) {eNB};

% Líneas de vida
\draw[dashed, black!30, thin] (0,-0.3) -- (0,-10.3);
\draw[dashed, black!30, thin] (2.5,-0.3) -- (2.5,-10.3);
\draw[dashed, black!30, thin] (8,-0.3) -- (8,-10.3);

% -----------------------------------------------------------------
% Paso 1: transmisión de MSG1
% -----------------------------------------------------------------

\node[step] at (-1.7,-0.85) {Paso 1};

\node[sidenote, anchor=north east] at (-0.3,-1.15)
    {Ambos UEs eligen\\el mismo preámbulo};

\draw[msg]
    (0,-1.0)
    --
    node[sloped, above, font=\small, pos=0.5]
    {MSG1: Preamble (RAPID $r$)}
    (8,-1.8);

\draw[msg]
    (2.5,-1.95)
    --
    (8,-2.5);

\node[callout, anchor=north west] at (8.3,-1.7)
    {\textbf{Evento en el eNB:}\\
     colisión no detectada};

% -----------------------------------------------------------------
% Paso 2: transmisión de MSG2
% -----------------------------------------------------------------

\node[step] at (-1.7,-3.55) {Paso 2};

\node[sidenote, anchor=north east] at (-0.3,-3.85)
    {El eNB envía un único\\RAR para el RAPID $r$};

\draw[msg]
    (8,-3.7)
    --
    node[sloped, above, font=\small, pos=0.5]
    {MSG2: RAR (un único TC-RNTI)}
    (0,-4.5);

\draw[msg]
    (8,-4.4)
    --
    (2.5,-4.95);

\node[sidenote, anchor=north west] at (8.3,-4.7)
    {Ambos UEs reciben\\el mismo RAR};

% -----------------------------------------------------------------
% Paso 3: transmisión simultánea de MSG3
% -----------------------------------------------------------------

\node[step] at (-1.7,-5.85) {Paso 3};

\node[sidenote, anchor=north east] at (-0.3,-6.15)
    {Ambos UEs transmiten\\en el mismo recurso};

\draw[msg]
    (0,-6.0)
    --
    node[sloped, above, font=\small, pos=0.5]
    {MSG3 sobre el mismo recurso}
    (8,-6.8);

\draw[msg]
    (2.5,-6.95)
    --
    (8,-7.5);

\node[callout, anchor=north west] at (8.3,-6.7)
    {\textbf{Evento en el eNB:}\\
     superposición de MSG3;\\
     ninguno puede decodificarse};

% -----------------------------------------------------------------
% Paso 4: fallo del procedimiento
% -----------------------------------------------------------------

\node[step] at (-1.7,-8.55) {Paso 4};

\node[sidenote, anchor=north east] at (-0.3,-8.85)
    {No se recibe MSG4};

\draw[
    ->,
    semithick,
    red!60!black,
    dashed
]
    (8,-8.6)
    --
    node[
        sloped,
        above,
        font=\small,
        text=red!60!black,
        pos=0.4
    ]
    {Sin respuesta}
    (0,-9.4);

\draw[
    ->,
    semithick,
    red!60!black,
    dashed
]
    (8,-9.3)
    --
    (2.5,-9.85);

\node[
    font=\small\bfseries,
    text=red!60!black
] at (4,-10.6)
    {Ambos UEs fallan --- backoff y reintento};

\end{tikzpicture}

\end{document}